\documentclass[preprint,a4paper,authoryear]{imsart}
%%%%%%%%%%%%%%%%% packages %%%%%%%%%%%%%%%%%%
%\usepackage{a4wide}
\usepackage{amssymb,amsthm,amsmath}
%\usepackage{bbm}
\usepackage[sort&compress]{natbib}
%\usepackage{cancel}
%\usepackage{mathrsfs}
\usepackage{hyperref}
%\usepackage{times}
%%%%%%%%%%%%%%%%% end packages %%%%%%%%%%%%%%%%%%


\numberwithin{equation}{section}

%%%%%%%%%%%%%%%% new commands %%%%%%%%%%%%%
\newcommand{\Cov}{\operatorname{Cov}}
\newcommand{\WN}{\operatorname{WN}}
\newcommand{\Var}{\operatorname{Var}}
\newcommand{\im}{\operatorname{im}}
\newcommand{\rank}{\operatorname{rank}}
\newcommand{\convd}{\xrightarrow{\mathcal{D}}}
\newcommand{\eqd}{\stackrel{\mathcal{D}}{=}}
\newcommand{\adj}{\operatorname{adj}}
\newcommand{\spa}{\operatorname{span}}
\newcommand{\re}{\operatorname{Re}}
\newcommand{\imag}{\operatorname{Im}}
\newcommand{\E}{\mathbb{E}}
\newcommand{\Pb}{\mathbb{P}}
\newcommand{\N}{\mathbb{N}}
\newcommand{\Z}{\mathbb{Z}}
\newcommand{\R}{\mathbb{R}}
\newcommand{\C}{\mathbb{C}}
\newcommand{\Y}{\boldsymbol{Y}}
\newcommand{\y}{\boldsymbol{y}}
\newcommand{\X}{\boldsymbol{X}}
\newcommand{\Lb}{\boldsymbol{L}}
\newcommand{\J}{\boldsymbol{J}}
\newcommand{\W}{\boldsymbol{W}}
\newcommand{\V}{\boldsymbol{V}}
\newcommand{\ZZ}{\boldsymbol{Z}}
\newcommand{\NN}{\boldsymbol{N}}
\newcommand{\M}{\boldsymbol{M}}
\newcommand{\beps}{\boldsymbol{\varepsilon}}
\newcommand{\bepsh}{\boldsymbol{\varepsilon}^{(h)}}
\newcommand{\bnu}{\boldsymbol{\nu}}
\newcommand{\bu}{\boldsymbol{u}}
\newcommand{\bx}{\boldsymbol{x}}
\newcommand{\bxi}{\boldsymbol{\xi}}
\newcommand{\bb}{\boldsymbol{b}}
\newcommand{\bs}{\boldsymbol{s}}
\newcommand{\bzero}{\boldsymbol{0}}
\newcommand{\be}{\boldsymbol{e}}
\newcommand{\vecM}{\mathscr{M}}
\newcommand{\nuL}{\nu^{\Lb}}
\newcommand{\hyp}{-}
%\newcommand{\bdoi}{\doi}
\newcommand{\burl}{\url}
%%%%%%%%%%%%%%%% end new commands %%%%%%%%%%%%%


\theoremstyle{plain}
\newtheorem{theorem}{Theorem}[section]
\newtheorem{proposition}[theorem]{Proposition}
\newtheorem{lemma}[theorem]{Lemma}
\newtheorem{corollary}[theorem]{Corollary}

\theoremstyle{definition}
\newtheorem{definition}[theorem]{Definition}

\theoremstyle{remark}
\newtheorem{remark}[theorem]{Remark}

\newcounter{assumptionlevy}
\newtheorem{assumptionlevy}[assumptionlevy]{Assumption}
\renewcommand{\theassumptionlevy}{{\bf L\arabic{assumptionlevy}}}

\newcounter{assumptioneigen}
\newtheorem{assumptioneigen}[assumptioneigen]{Assumption}
\renewcommand{\theassumptioneigen}{{\bf E\arabic{assumptioneigen}}}


\begin{document}

 
\begin{frontmatter}

\title{Correction to: Multivariate CARMA processes, continuous-time state space models and complete regularity of the innovations of the sampled processes, Bernoulli 18, pp. 46-63, 2012}%\protect\thanksref{T1}%}
\runtitle{}
%\thankstext{T1}{Footnote to the title with the `thankstext' command.}
\begin{aug}
	%%%%%%%%%%%%%%%%%%%%%%%%%%%%%%%%%%%%%%%%%%%%%%%
	%% ORCID can be inserted by command:         %%
	%% \orcid{0000-0000-0000-0000}               %%
	%%%%%%%%%%%%%%%%%%%%%%%%%%%%%%%%%%%%%%%%%%%%%%%
	\author[A]{\fnms{Robert}~\snm{Stelzer}\ead[label=e1]{robert.stelzer@uni-ulm.de}\ead[label=u1,url]{https://www.uni-ulm.de/mawi/finmath}}
%	\orcid{0000-0002-2700-0735}

	%%%%%%%%%%%%%%%%%%%%%%%%%%%%%%%%%%%%%%%%%%%%%%
	%% Addresses                                %%
	%%%%%%%%%%%%%%%%%%%%%%%%%%%%%%%%%%%%%%%%%%%%%%
	\address[A]{Institute of Mathematical Finance, Ulm University, Helmholtzstrasse 18, D-89081 Ulm, Germany\printead[presep={,\ }]{e1,u1}}
	
\end{aug}


\begin{abstract}
A serious flaw in the proof of the equivalence of continuous time state space models and MCARMA processes spotted in \cite{FasenSchenk2024} is corrected. We point out that likewise an issue in the proof of Theorem 3.2 in \cite{brockwell2011parametric} can be resolved and, hence, any MCARMA process and linear state space model has both a controller and an observer canonical representation. Equivalently, the transfer function has both a left and right matrix fraction representation.
\end{abstract}

\begin{keyword}[class=AMS]
\kwd[Primary ]{60G10} %% Stationary Processes
\kwd[; secondary ]{60G51} %% Processes with independent increments
\end{keyword}

\begin{keyword}
\kwd{multivariate CARMA process}
\kwd{state space representation}
\end{keyword}

\end{frontmatter}
\section{Introduction}
Multivariate CARMA processes as introduced in \cite{marquardt2007multivariate} are used in various applications and have also been implemented in R packages (see e.g. \cite{Tomasson2018}).
The statistical inference theory developed in \cite{FasenKimmig2020,FasenKimmig2017,FasenMayer2022,FasenScholz2019,SchlemmStelzer2012EJS}, for instance, hinges crucially on the equivalence of the class of L\'evy-driven MCARMA processes to the class  of L\'evy-driven linear state space models, as identifiability is typically ensured by considering state space models in echolon form.

But as observed in \cite{FasenSchenk2024} the proof of this equivalence (Corollary 3.4) in \cite{schlemmmixing2010} is incorrect, since Appendix 2 of \cite{Caines1988} does not guarantee that the leading coefficient of the ``denominator'' in the left matrix fraction representation of the transfer function can be chosen to be the identity. The same problem - likewise spotted by \cite{FasenSchenk2024} -  with right matrix fractions arises in the proof of Theorem 3.2 in \cite{brockwell2011parametric}, where the authors refer to Lemma 6.3-8 of \cite{Kailath1980}, which also does not establish that the leading coefficient in the ``denominator'' can be taken to be the identity. 

Hence, we shall prove below that any L\'evy-driven linear state space model can be represented in both observer and controller canonical form and that its transfer function always has both a left and a right matrix fraction representation with the ``denominator'' having the identity as the leading coefficient. This in particular implies that Corollary 3.4 in \cite{schlemmmixing2010} and Theorem 3.2 in \cite{brockwell2011parametric} are correct.

In the following we refer the reader to our original paper \cite{schlemmmixing2010} for any unexplained notions and notation. We will refer to equations, definitions, theorems etc. in that paper by putting the letter ``P'' in front of the respective number.

\section{Results}
Note first that there is a typo in Equation (P3.4b). In line with  \cite[Theorem 3.12]{marquardt2007multivariate}  it should read
\[
\beta_{p-j} = I_{\{0,\ldots,q\}}(j)\left[-\sum_{i=1}^{p-j-1}{A_i\beta_{p-j-i}+B_{q-j}}\right].\qquad\qquad\qquad\qquad\hfill \mathrm{(P3.4b)}
\]
\begin{proposition}[Observer canonical form/MCARMA process representation]\label{prop:observer}
	For the output process $\mathbf{Y}$ of a L\'evy-driven state space model of the form (P3.5) there exist a $p\in\mathbb{N}$ and matrices $\mathcal{A,B,C}$ of the form (P3.4) such that $\mathbf{Y}$  is the output process of an MCARMA state space representation (linear state space model in observer canonical form) of the form (P3.3).
\end{proposition}
\begin{proof}
	Denote by $H(z)=C(z\mathbb{I}_N-A)^{-1} B$ the transfer function of the linear state space model (P3.5). Two linear state space models driven by the same L\'evy process produce the same output process $\mathbf{Y}$ if they have the same transfer function (see e.g. Lemma 3.2 of \cite{SchlemmStelzer2012EJS}). 
	
	Necessarily, every element of the matrix $H(z)$  is a rational function in $z$ with the degree of the denominator exceeding the degree of the numerator (see p. 106 in \cite{Brockett2015}). Now Proof 2 of Theorem 17.1 of \cite{Brockett2015} gives matrices $\mathcal{A,B,C}$ of the form (P3.4) with $H(z)=\mathcal{C}(z\mathbb{I}_N-\mathcal{A})^{-1} \mathcal{B}$. These matrices define a model of the  form (P3.3) with the given output process $\mathbf Y$.
	
\end{proof}
\begin{proposition}[Controller canonical form]\label{prop:controller}
	For the output process $\mathbf{Y}$ of a L\'evy-driven state space model of the form (P3.5) there exist a $p\in\mathbb{N}$ and matrices $\mathfrak{A,B,C}$ of the form 
	\begin{align}
	 \mathfrak{A} =& \left(\begin{array}{ccccc}
			0 & \mathbb{I}_m & 0 & \ldots & 0 \\
			0 & 0 & \mathbb{I}_m & \ddots & \vdots \\
			\vdots && \ddots & \ddots & 0\\
			0 & \ldots & \ldots & 0 & \mathbb{I}_m\\
			-\widetilde A_p & -\widetilde A_{p-1} & \ldots & \ldots & -\widetilde A_1
		\end{array}\right)\in M_{pm}(\R),\\
	\mathfrak{B}=&\left(0_m,\ldots,0_m,\mathbb{I}_m\right)^T\in M_{pm,m}(\R)\text{ and}\\
	\mathfrak{C}=&\left(\begin{array}{ccc}\widetilde B_0, & \ldots, & \widetilde B_{p-1}\end{array}\right)\in M_{d,pm}(\mathbb{R})
	\end{align}
such that $\mathbf{Y}$  is the output process of  a linear state space model in controller canonical form 
	\begin{equation}
		\mathrm{d}\widetilde{\boldsymbol{G}}(t)=\mathfrak{A}\widetilde{\boldsymbol{G}}(t)\mathrm{d}t+\mathfrak{B}\mathrm{d}\Lb(t),\quad \Y(t)=\mathfrak{C}\widetilde{\boldsymbol{G}}(t),\quad t\in\R.
	\end{equation}
	.
\end{proposition}
\begin{proof}
The proof is analogous to the one of Proposition \ref{prop:controller} only that one now uses Proof 1 of Theorem 17.1 of \cite{Brockett2015}.
\end{proof}
\begin{proposition}[Matrix fraction representations]
	Let $H$ be the transfer function  of a L\'evy-driven  state space model of the form (P3.5). There exist $p,q,\tilde q\in\mathbb{N}_0$, $p>q, p>\tilde q$ and
	polynomials $P\in M_d(\mathbb{R}[z]),Q\in M_{d,m}(\mathbb{R}[z])$ (left matrix fraction description) as well as $\widetilde P \in M_m(\mathbb{R}[z]), \widetilde Q\in M_{d,m}(\mathbb{R}[z])$ (right matrix fraction description) such that
	\[
	H(z)=C(z\mathbb{I}_N-A)^{-1} B=P^{-1}(z)Q(z)=\widetilde Q(z) \widetilde P^{-1}(z)\,\,\forall z \in \mathbb{C}.
	\]
	with
	\begin{eqnarray*}
		P(z)=\mathbb{I}_dz^p+A_1z^{p-1}+\ldots+A_p, &\qquad
		Q(z)=B_0z^q+B_1 z^{q-1}+\ldots+B_q,\\
	\widetilde	P(z)=\mathbb{I}_mz^p+\widetilde A_1z^{p-1}+\ldots+\widetilde A_p,&\qquad
\widetilde	Q(z)=\widetilde B_0z^{\tilde q}+\widetilde B_1 z^{\tilde q-1}+\ldots+\widetilde B_{\tilde q}.
	\end{eqnarray*}
\end{proposition}
\begin{proof}
	\emph{Left matrix fraction:}
Let  $\mathcal{A,B,C}$ be the matrices obtained by Proposition \ref{prop:observer}, define $P,p$ by the elements/dimensions of $\mathcal A$, set $q=p- \min(\{i=1,\ldots,p:\,\beta_i\not =0\}$ and $B_{q-j}=\beta_{p-j}+\sum_{i=1}^{p-j-1}{A_i\beta_{p-j-i}},\,j=0,\ldots,q$. Defining $Q$ accordingly the arguments in the first step of the proof of Theorem P3.3 combined with   \cite{marquardt2007multivariate},  Theorem 3.12, establish that $H(z)=P^{-1}(z)Q(z)$.

	\emph{Right matrix fraction:}
Likewise let  $\mathfrak{A,B,C}$  be the matrices obtained by Proposition \ref{prop:controller}  and define $\widetilde P, \widetilde Q, p, \tilde q$  via their elements/dimensions. Then the arguments in the proof of  \cite{brockwell2011parametric}, Theorem 3.2, show 	$H(z)=\widetilde Q(z) \widetilde P^{-1}(z)$.

That the degrees of the polynomials $P,\widetilde P$ agree follows from comparing Proof 1 and 2 of  Theorem 17.1 of \cite{Brockett2015}.
\end{proof}
Whereas in the original paper and in \cite{brockwell2011parametric}, respectively, the existence of the observer and controller canonical form, respectively, was incorrectly  a consequence of the existence of the matrix left and right factorization with the leading coefficient of the ``denominator'' being the identity, in our above flow of arguments it is now actually the other way round.


\bibliography{library}
\bibliographystyle{imsart-nameyear}

\end{document}

\begin{thebibliography}{13}
% BibTex style file: imsart-nameyear.bst, 2017-11-03
% Default style options (sort=1,type=nameyear).
% Used options (sort=1,type=nameyear).

\bibitem[\protect\citeauthoryear{Brockett}{2015}]{Brockett2015}
\begin{bbook}[author]
\bauthor{\bsnm{Brockett},~\bfnm{Roger~W.}\binits{R.~W.}}
(\byear{2015}).
\btitle{Finite dimensional linear systems}.
\bseries{Classics in Applied Mathematics}
\bvolume{74}.
\bpublisher{Society for Industrial and Applied Mathematics (SIAM),
  Philadelphia, PA}
\bnote{Reprint of the 1970 original}.
\bdoi{10.1137/1.9781611973884}
\bmrnumber{3486166}
\end{bbook}
\endbibitem

\bibitem[\protect\citeauthoryear{Brockwell and
  Schlemm}{2013}]{brockwell2011parametric}
\begin{barticle}[author]
\bauthor{\bsnm{Brockwell},~\bfnm{Peter~J.}\binits{P.~J.}} \AND
  \bauthor{\bsnm{Schlemm},~\bfnm{Eckhard}\binits{E.}}
(\byear{2013}).
\btitle{Parametric estimation of the driving {L}\'evy process of multivariate
  {CARMA} processes from discrete observations}.
\bjournal{J. Multivariate Anal.}
\bvolume{115}
\bpages{217--251}.
\bdoi{10.1016/j.jmva.2012.09.004}
\bmrnumber{3004556}
\end{barticle}
\endbibitem

\bibitem[\protect\citeauthoryear{Caines}{2018}]{Caines1988}
\begin{bbook}[author]
\bauthor{\bsnm{Caines},~\bfnm{Peter~E.}\binits{P.~E.}}
(\byear{2018}).
\btitle{Linear stochastic systems}.
\bseries{Classics in Applied Mathematics}
\bvolume{77}.
\bpublisher{Society for Industrial and Applied Mathematics (SIAM),
  Philadelphia, PA}
\bnote{Reprint of the 1988 original [MR0944080]}.
\bmrnumber{3908672}
\end{bbook}
\endbibitem

\bibitem[\protect\citeauthoryear{Fasen and Kimmig}{2017}]{FasenKimmig2017}
\begin{barticle}[author]
\bauthor{\bsnm{Fasen},~\bfnm{Vicky}\binits{V.}} \AND
  \bauthor{\bsnm{Kimmig},~\bfnm{Sebastian}\binits{S.}}
(\byear{2017}).
\btitle{Information criteria for multivariate {CARMA} processes}.
\bjournal{Bernoulli}
\bvolume{23}
\bpages{2860--2886}.
\bdoi{10.3150/16-BEJ830}
\bmrnumber{3648048}
\end{barticle}
\endbibitem

\bibitem[\protect\citeauthoryear{Fasen-Hartmann and
  Kimmig}{2020}]{FasenKimmig2020}
\begin{barticle}[author]
\bauthor{\bsnm{Fasen-Hartmann},~\bfnm{Vicky}\binits{V.}} \AND
  \bauthor{\bsnm{Kimmig},~\bfnm{Sebastian}\binits{S.}}
(\byear{2020}).
\btitle{Robust estimation of stationary continuous-time {ARMA} models via
  indirect inference}.
\bjournal{J. Time Series Anal.}
\bvolume{41}
\bpages{620--651}.
\bdoi{10.1111/jtsa.12526}
\bmrnumber{4176167}
\end{barticle}
\endbibitem

\bibitem[\protect\citeauthoryear{Fasen-Hartmann and
  Mayer}{2022}]{FasenMayer2022}
\begin{barticle}[author]
\bauthor{\bsnm{Fasen-Hartmann},~\bfnm{Vicky}\binits{V.}} \AND
  \bauthor{\bsnm{Mayer},~\bfnm{Celeste}\binits{C.}}
(\byear{2022}).
\btitle{Whittle estimation for continuous-time stationary state space models
  with finite second moments}.
\bjournal{Ann. Inst. Statist. Math.}
\bvolume{74}
\bpages{233--270}.
\bdoi{10.1007/s10463-021-00802-6}
\bmrnumber{4396674}
\end{barticle}
\endbibitem

\bibitem[\protect\citeauthoryear{Fasen-Hartmann and
  Schenk}{2024}]{FasenSchenk2024}
\begin{barticle}[author]
\bauthor{\bsnm{Fasen-Hartmann},~\bfnm{Vicky}\binits{V.}} \AND
  \bauthor{\bsnm{Schenk},~\bfnm{Lea}\binits{L.}}
(\byear{2024}).
\btitle{Mixed orthogonality graphs for continuous-time state space models and
  orthogonal projections}.
\bjournal{Journal of Time Series Analysis}.
\bnote{to appear}.
\bdoi{https://doi.org/10.1111/jtsa.12787}
\end{barticle}
\endbibitem

\bibitem[\protect\citeauthoryear{Fasen-Hartmann and
  Scholz}{2019}]{FasenScholz2019}
\begin{barticle}[author]
\bauthor{\bsnm{Fasen-Hartmann},~\bfnm{Vicky}\binits{V.}} \AND
  \bauthor{\bsnm{Scholz},~\bfnm{Markus}\binits{M.}}
(\byear{2019}).
\btitle{Quasi-maximum likelihood estimation for cointegrated continuous-time
  linear state space models observed at low frequencies}.
\bjournal{Electron. J. Stat.}
\bvolume{13}
\bpages{5151--5212}.
\bdoi{10.1214/19-EJS1636}
\bmrnumber{4042406}
\end{barticle}
\endbibitem

\bibitem[\protect\citeauthoryear{Kailath}{1980}]{Kailath1980}
\begin{bbook}[author]
\bauthor{\bsnm{Kailath},~\bfnm{Thomas}\binits{T.}}
(\byear{1980}).
\btitle{Linear systems}.
\bseries{Prentice-Hall Information and System Sciences Series}.
\bpublisher{Prentice-Hall, Inc., Englewood Cliffs, NJ}.
\bmrnumber{569473}
\end{bbook}
\endbibitem

\bibitem[\protect\citeauthoryear{Marquardt and
  Stelzer}{2007}]{marquardt2007multivariate}
\begin{barticle}[author]
\bauthor{\bsnm{Marquardt},~\bfnm{Tina}\binits{T.}} \AND
  \bauthor{\bsnm{Stelzer},~\bfnm{Robert}\binits{R.}}
(\byear{2007}).
\btitle{Multivariate {CARMA} processes}.
\bjournal{Stochastic Process. Appl.}
\bvolume{117}
\bpages{96--120}.
\bdoi{10.1016/j.spa.2006.05.014}
\bmrnumber{2287105}
\end{barticle}
\endbibitem

\bibitem[\protect\citeauthoryear{Schlemm and
  Stelzer}{2012a}]{SchlemmStelzer2012EJS}
\begin{barticle}[author]
\bauthor{\bsnm{Schlemm},~\bfnm{Eckhard}\binits{E.}} \AND
  \bauthor{\bsnm{Stelzer},~\bfnm{Robert}\binits{R.}}
(\byear{2012}a).
\btitle{Quasi maximum likelihood estimation for strongly mixing state space
  models and multivariate {L}\'evy-driven {CARMA} processes}.
\bjournal{Electron. J. Stat.}
\bvolume{6}
\bpages{2185--2234}.
\bdoi{10.1214/12-EJS743}
\bmrnumber{3020261}
\end{barticle}
\endbibitem

\bibitem[\protect\citeauthoryear{Schlemm and
  Stelzer}{2012b}]{schlemmmixing2010}
\begin{barticle}[author]
\bauthor{\bsnm{Schlemm},~\bfnm{Eckhard}\binits{E.}} \AND
  \bauthor{\bsnm{Stelzer},~\bfnm{Robert}\binits{R.}}
(\byear{2012}b).
\btitle{Multivariate {CARMA} processes, continuous-time state space models and
  complete regularity of the innovations of the sampled processes}.
\bjournal{Bernoulli}
\bvolume{18}
\bpages{46--63}.
\bdoi{10.3150/10-BEJ329}
\bmrnumber{2888698}
\end{barticle}
\endbibitem

\bibitem[\protect\citeauthoryear{T{\'o}masson}{2018}]{Tomasson2018}
\begin{binproceedings}[author]
\bauthor{\bsnm{T{\'o}masson},~\bfnm{Helgi}\binits{H.}}
(\byear{2018}).
\btitle{Implementation of Multivariate Continuous-time ARMA Models}.
In \bbooktitle{Continous-Time Modeling in the Behavioral and Related Sciences}
\bpages{359-387}.
\bpublisher{Springer}, \baddress{Cham}.
\end{binproceedings}
\endbibitem

\end{thebibliography}
